% ==============================================================================
% Scenario 3: Fault Tolerance and Blast Radius Containment
% Manuscript Results Section for Section VI
% Active Enforcement Mode (cic-v5.3) - 30 Measured Benchmark Trials
% ==============================================================================

\subsection{Scenario 3: Fault Tolerance and Blast Radius Containment}
\label{sec:scenario3_fault_tolerance}

To evaluate the operational resilience of autonomous DDoS mitigation under control-plane component failures, Scenario~3 compares the fault isolation capabilities of the decentralized architecture against the monolithic centralized coordinator. In microservice environments, a monitoring or mitigation failure within one service's control loop must not compromise unaffected services or induce cascade failures across the service mesh.

\subsubsection{Experimental Design and Failure Injection Methodology}
We conduct two complementary failure injection experiments across thirty measured trials ($N=10$ per condition) on the Istio-enabled Bookinfo benchmark, enforcing strict 1:1 hardware resource parity ($4 \times 1000\,\text{m}$ CPU, $4 \times 2048\,\text{Mi}$ RAM decentralized vs.\ $1 \times 4000\,\text{m}$ CPU, $1 \times 8192\,\text{Mi}$ RAM centralized):

\begin{enumerate}
    \item \textbf{Experiment A (Transient Defender Failure):} Simulates an abrupt controller crash by deleting the \texttt{reviews} defender pod (\texttt{replicas=1}). We track the full recovery timeline: pod scheduling and recreation, earliest internal thread initialization in container logs, Kubelet readiness probe passage (\texttt{Ready=True}), and external status confirmation of active scraping and detection loops. Sibling defender continuity and Kubernetes Lease lock integrity are continuously audited.
    \item \textbf{Experiment B (Sustained Outage and Blast Radius):} Evaluates multi-service defense under sustained controller unavailability during a concurrent $125\,\text{RPS}$ DDoS attack ($65\,\text{RPS}$ \texttt{productpage}, $20\,\text{RPS}$ \texttt{details}, $20\,\text{RPS}$ \texttt{reviews}, $20\,\text{RPS}$ \texttt{ratings}) and $24\,\text{RPS}$ legitimate mesh traffic. In the decentralized architecture (Experiment~B1), the \texttt{reviews} defender is scaled to zero replicas, measuring whether surviving sibling controllers continue autonomous detection and mitigation. In the centralized baseline (Experiment~B2), the monolithic coordinator is scaled to zero replicas without automated failover, deliberately establishing the single-point-of-failure (SPOF) blast radius under identical hardware resources.
\end{enumerate}

Every trial enforces a mandatory 15-second inter-trial settling barrier with automated baseline route and Lease restoration.

\subsubsection{Transient Recovery and Self-Healing Dynamics}
Table~\ref{tab:scenario3_recovery_lifecycle} summarizes the empirical lifecycle milestones observed during transient pod failures across ten independent trials. All statistics reflect sample standard deviations ($s$, $N-1=9$) and nearest-observation percentiles.

\begin{table}[!t]
\renewcommand{\arraystretch}{1.2}
\caption{Experiment A: Transient Defender Failure and Self-Healing Milestones ($N=10$ Trials).}
\label{tab:scenario3_recovery_lifecycle}
\centering
\footnotesize
\setlength{\tabcolsep}{4.0pt}
\begin{tabular}{@{}l c c c c c@{}}
\toprule
\textbf{Lifecycle Milestone} & \textbf{Mean (s)} & \textbf{Std Dev (s)} & \textbf{Min (s)} & \textbf{Max (s)} & \textbf{p95 (s)} \\
\midrule
Replacement Pod Created     & 2.075 & 0.806 & 1.557 & 4.304 & 4.304 \\
Internal Thread Startup     & 4.277 & 2.670 & 3.167 & 11.856 & 11.856 \\
Readiness Probe (1/1)       & 12.638 & 4.261 & 10.414 & 20.957 & 20.957 \\
External Loop Confirmation  & 13.193 & 4.266 & 10.965 & 21.537 & 21.537 \\
\midrule
\textbf{Total Observed Downtime} & \textbf{13.193} & \textbf{4.266} & \textbf{10.965} & \textbf{21.537} & \textbf{21.537} \\
\bottomrule
\end{tabular}
\end{table}

Following pod deletion, the Kubernetes ReplicaSet controller detected termination and scheduled a replacement pod within $2.075 \pm 0.806\,\text{s}$ (peaking at $4.304\,\text{s}$ in Trial~1). Container startup logs reveal that internal background scraping and detection threads initialized at $4.277 \pm 2.670\,\text{s}$. The Kubelet HTTP readiness probe confirmed endpoint readiness at $12.638 \pm 4.261\,\text{s}$, while external status polling verified active monitoring thread resumption at $13.193 \pm 4.266\,\text{s}$.

Throughout all ten transient failure trials, continuous inspection confirmed that sibling defenders (\texttt{productpage}, \texttt{details}, \texttt{ratings}) maintained a \textbf{100.0\% continuity pass rate} ($30/30$ controller observations), experiencing zero container restarts, zero dropped Lease locks, and continuous thread execution. Furthermore, distributed Lease auditing confirmed that the \texttt{reviews-mitigation-lock} remained unheld without dangling lockouts or split-brain deadlocks.

\subsubsection{Blast Radius Containment Under Sustained Outage}
Table~\ref{tab:scenario3_blast_radius} presents the measured blast radius and attacker diversion metrics under sustained controller failure during concurrent multi-service DDoS attacks.

\begin{table*}[!t]
\renewcommand{\arraystretch}{1.2}
\caption{Experiment B: Fault Isolation and Attacker Diversion Under Sustained Outage ($N=20$ Trials).}
\label{tab:scenario3_blast_radius}
\centering
\footnotesize
\setlength{\tabcolsep}{5pt}
\begin{tabular}{@{}l l l c c c c c c@{}}
\toprule
\textbf{Architecture} & \textbf{Outage Condition} & \textbf{Service} & \textbf{Controller} & \textbf{Quarantine Success} & \textbf{Decoy Hits (Att.)} & \textbf{Benign Leakage} & \textbf{Retained Defense} & \textbf{Blast Radius} \\
\midrule
\textbf{Decentralized} & \texttt{reviews} Defender & \texttt{productpage} & Online  & 10/10 (100\%) & 18,409\textsuperscript{*} & 0 & 100\% & 0\% \\
(Experiment B1)        & Scaled to 0               & \texttt{details}     & Online  & 10/10 (100\%) & 5,318                     & 0 & 100\% & 0\% \\
                       &                           & \texttt{reviews}     & Offline & 0/10 (0\%)    & 0                         & 0 & 0\%   & \textbf{100\%} \\
                       &                           & \texttt{ratings}     & Online  & 10/10 (100\%) & 4,679                     & 0 & 100\% & 0\% \\
\cmidrule(l){3-9}
                       & \textbf{Summary}          & —                    & 3 of 4  & \textbf{30/30 (100\%)} & \textbf{28,406 Captured} & \textbf{0 (0.0\%)} & \textbf{75.0\%} & \textbf{25.0\%} \\
\midrule
\textbf{Centralized}   & Coordinator               & \texttt{productpage} & Offline & 0/10 (0\%)    & 0                         & 0 & 0\%   & 100\% \\
(Experiment B2)        & Scaled to 0               & \texttt{details}     & Offline & 0/10 (0\%)    & 0                         & 0 & 0\%   & 100\% \\
                       &                           & \texttt{reviews}     & Offline & 0/10 (0\%)    & 0                         & 0 & 0\%   & 100\% \\
                       &                           & \texttt{ratings}     & Offline & 0/10 (0\%)    & 0                         & 0 & 0\%   & 100\% \\
\cmidrule(l){3-9}
                       & \textbf{Summary}          & —                    & 0 of 1  & \textbf{0/40 (0\%)}    & \textbf{0 Captured}      & \textbf{0 (0.0\%)} & \textbf{0.0\%}  & \textbf{100.0\%} \\
\bottomrule
\multicolumn{9}{@{}l}{\footnotesize \textsuperscript{*}Positive diversion confirmed via request-ID matching; aggregate productpage count labeled as unresolved due to generator drain buffering.}
\end{tabular}
\end{table*}

The empirical data demonstrates a fundamental structural divergence in failure containment:
\begin{itemize}
    \item \textbf{Decentralized Architecture (25\% Blast Radius):} In all ten trials ($30/30$ surviving service instances), the \texttt{productpage}, \texttt{details}, and \texttt{ratings} controllers operated completely unaffected by the \texttt{reviews} outage. Each surviving defender autonomously executed its detection loop, acquired its Lease, and patched its respective VirtualService, capturing a cumulative $28,406$ attacker requests across the three decoys. Only the unprotected \texttt{reviews} service experienced unmitigated attack traffic. The failure blast radius was strictly bounded to \textbf{25.0\%}, preserving \textbf{75.0\%} of active defense mesh-wide.
    \item \textbf{Centralized Architecture (100\% Blast Radius):} Because the centralized coordinator represents a single control plane, its outage simultaneously dismantled defense across all four microservices ($0/40$ service instances quarantined, $0$ decoy hits). The failure blast radius expanded to \textbf{100.0\%}, causing total protection collapse ($0.0\%$ retained defense).
    \item \textbf{Decoy Isolation Integrity:} Across all trials in both architectures, exactly zero benign requests were diverted to decoys ($0.0\%$ leakage), verifying that quarantine rules strictly targeted attacker identities.
\end{itemize}

\subsubsection{Impact on Legitimate Mesh Traffic Quality of Service}
Table~\ref{tab:scenario3_qos_comparison} analyzes background legitimate HTTP traffic performance ($N=14,400$ total requests across 20 trials).

\begin{table}[!t]
\renewcommand{\arraystretch}{1.2}
\caption{Legitimate Traffic Performance During Sustained Defender Outages ($N=20$ Trials).}
\label{tab:scenario3_qos_comparison}
\centering
\footnotesize
\setlength{\tabcolsep}{3.5pt}
\begin{tabular}{@{}l c c c c@{}}
\toprule
\textbf{Metric} & \textbf{Decentralized} & \textbf{Centralized} & \textbf{Difference} & \textbf{Significance} \\
\midrule
Requests Evaluated & 7,200 & 7,200 & — & — \\
Success Rate (\%)  & 99.74\% & 99.71\% & $+0.03\%$ & $p = 0.875$ \\
Errors per Trial   & $1.90 \pm 2.51$ & $2.10 \pm 2.42$ & $-9.5\%$ & $p = 0.858$ \\
Mean Latency       & $17.13 \pm 9.00$\,ms & $49.53 \pm 16.84$\,ms & \textbf{$-65.4\%$} & $\mathbf{p = 1.055 \times 10^{-4}}$ \\
Median (p50)       & $4.49$\,ms ($s=0.81$) & $5.14$\,ms ($s=1.40$) & $-12.6\%$ & $p = 0.2215$ \\
Tail Latency (p95) & $45.18 \pm 28.28$\,ms & $149.47 \pm 49.66$\,ms & \textbf{$-69.8\%$} & $\mathbf{p = 4.477 \times 10^{-5}}$ \\
Extreme (p99)      & $247.56$\,ms ($s=277.34$) & $723.92$\,ms ($s=289.19$) & \textbf{$-65.8\%$} & $\mathbf{p = 0.00144}$ \\
\bottomrule
\end{tabular}
\end{table}

When the centralized coordinator fails, unmitigated attack traffic across all four services saturates Envoy sidecar connection pools, degrading service mesh performance. Legitimate mean latency increased from $17.13\,\text{ms}$ under decentralized containment to $49.53\,\text{ms}$ under centralized failure ($t = -5.3667, p = 1.055 \times 10^{-4}$ via Welch's two-sided t-test). More critically, tail p95 latency experienced a \textbf{3.31-fold degradation}, increasing from $45.18\,\text{ms}$ to $149.47\,\text{ms}$ ($t = -5.7710, p = 4.477 \times 10^{-5}$). Both QoS degradations are statistically significant at $\alpha = 0.001$, proving that the decentralized architecture's ability to contain defender failures directly preserves legitimate user quality of service. Extreme tail latency (p99) similarly showed statistically significant degradation ($247.56\,\text{ms}$ vs.\ $723.92\,\text{ms}$, $t = -3.7595, p = 0.00144$). Conversely, median latency (p50) differences ($4.49\,\text{ms}$ vs.\ $5.14\,\text{ms}$) did not reach statistical significance ($t = -1.2780, p = 0.2215$), confirming that unblocked requests experience normal baseline propagation.

\subsubsection{Scientific Attribution and Telemetry Nuances}
Three important experimental nuances warrant explicit scientific documentation:
\begin{enumerate}
    \item \textbf{Decoy Attacker Diversion Verification:} Attacker diversion was empirically established through request-ID matching: samples of active attack tokens (\texttt{bench-att-s3-b1-t*}) were verified within decoy FIFO memory buffers for all surviving services, while zero matching tokens reached the decoy for \texttt{reviews}. On \texttt{productpage}, the cumulative counter increase ($18,409$) exceeded the nominal steady-state formula expectation ($65\,\text{RPS} \times (30 - 5.97\,\text{s}) \approx 1,562\,\text{reqs}$). This discrepancy is attributed to initial process connection overhead and multi-threaded generator queue buffering that drained through the Istio Ingress Gateway after route insertion. Following rigorous reporting practices, positive attacker diversion is confirmed via request-ID matching, while the cumulative aggregate count is conservatively noted as unresolved.
    \item \textbf{Data Recording Methodology and Scope:} To prevent command-pipe saturation and out-of-memory crashes on the Minikube virtual machine runner during 30 consecutive trials, the harness recorded per-trial aggregate distribution statistics ($N=720$ requests per trial, $14,400$ total) rather than serializing individual per-request transaction traces. All reported confidence intervals and significance tests are computed over these trial-level distributions.
    \item \textbf{Coordination Safety Across 10 Leases:} Audit verification of all ten Kubernetes Lease objects confirmed zero lock starvation or dangling ownership throughout all 30 trials. Distributed fencing tokens incremented monotonically, verifying that individual defender failures do not compromise distributed consensus across surviving controllers.
\end{enumerate}
